\section{支撑材料文件说明}

计算所用程序和结果数据保存在\path{supporting_materials}目录。限于篇幅，本附录选录与各模型直接对应的核心算法，主要文件见表\ref{tab:appendix-files}。

\begin{table}[H]
\centering
\caption{主要程序及功能}
\label{tab:appendix-files}
\begin{tabularx}{\textwidth}{p{0.43\textwidth}X}
\toprule
文件 & 功能 \\
\midrule
\path{code/analyze_p1.py} & 构造申请事件，分析需求、项目工时、科室结构与医生容量 \\
\path{code/build_unified_inputs.py} & 生成统一事件表、项目表和项目设备兼容集合 \\
\path{code/run_unified_scheduler.py} & 患者级可行性内核及问题二纯住院调度 \\
\path{code/run_pre_freeze_scenarios.py} & 问题二候选策略与问题三背景下限扫描 \\
\path{code/benchmark_exact_scheduler.py} & 问题二CP-SAT精确子问题对照 \\
\path{code/benchmark_exact_p3_scheduler.py} & 问题三联合调度CP-SAT精确子问题对照 \\
\path{code/verify_unified_schedule.py} & 从导出排程独立重建并检查全部约束 \\
\path{results/p1/} & 问题一需求、工时、峰平阈值与预测检验结果 \\
\path{results/final_frozen/} & 问题二、问题三及敏感性结果的机器可读文件 \\
\bottomrule
\end{tabularx}
\end{table}

计算顺序与正文一致：先完成数据预处理和问题一分析，再构建统一调度输入并求解问题二、三，随后计算精确子问题、校准和敏感性情景，最后独立检验排程可行性。各程序从题目原始附件或上一步结果读取数据，完整字段说明保存在同目录说明文档中。

\section{数据预处理程序}

数据预处理先由\path{analyze_p1.py}中的\path{expand_embedded_project_lists}按源系统明确的\path{][}边界拆分项目，再由\path{build_unified_inputs.py}中的\path{load_source}按“来源、患者编号、精确开单时刻”聚合事件、删除精确重复明细并生成稳定事件编号。下面直接摘录实际源码；函数后续的项目表字段整理保留在完整文件中。

\lstinputlisting[style=paperpython,firstline=218,lastline=238,firstnumber=218,caption={显式复合项目拆分的实际源码}]{../supporting_materials/code/analyze_p1.py}

\lstinputlisting[style=paperpython,firstline=141,lastline=177,firstnumber=141,caption={事件聚合与去重的实际源码}]{../supporting_materials/code/build_unified_inputs.py}

\lstinputlisting[style=paperpython,firstline=179,lastline=201,firstnumber=179,caption={模拟边界与稳定事件编号的实际源码}]{../supporting_materials/code/build_unified_inputs.py}

\section{问题一程序}

问题一按来源构造完整日序列，严格以2024年4月1日为时间分界，在训练期估计候选预测器并在留出期计算误差。对数岭回归由线性方程直接求解，峰平阈值只读取训练子集。以下均为\path{analyze_p1.py}的实际代码。

\lstinputlisting[style=paperpython,firstline=533,lastline=590,firstnumber=533,caption={时间留出预测与误差计算的实际源码}]{../supporting_materials/code/analyze_p1.py}

\lstinputlisting[style=paperpython,firstline=593,lastline=615,firstnumber=593,caption={训练期峰平分位数的实际源码}]{../supporting_materials/code/analyze_p1.py}

\section{问题二程序}

问题二由\path{task_order_candidates}生成四类项目顺序，\path{tentative_plan}逐项预留日历；任一任务失败时逆序释放此前位置，全部成功时保留预留结果。\path{plan_event}先尝试共同设备，再放开跨室安排。下面直接摘录\path{run_unified_scheduler.py}中的实现。

\lstinputlisting[style=paperpython,firstline=484,lastline=531,firstnumber=484,caption={候选项目序列与失败回滚的实际源码}]{../supporting_materials/code/run_unified_scheduler.py}

\lstinputlisting[style=paperpython,firstline=534,lastline=559,firstnumber=534,caption={共同设备优先的患者级排程实际源码}]{../supporting_materials/code/run_unified_scheduler.py}

\section{问题三程序}

问题三的正式下限按全评价期背景事件总数计算。\path{build_p3_candidate}先生成住院优先基线；若其未达到下限，则从空日历执行\path{global_deficit_repair}分支，先排至全期背景目标，再排住院和剩余背景。\path{solve_p3_epsilon}仅从达到下限的候选中按住院完成数及等待指标选择。下面直接摘录\path{run_pre_freeze_scenarios.py}中的两个实际函数。

\lstinputlisting[style=paperpython,firstline=216,lastline=293,firstnumber=216,caption={住院基线与全局缺口重构分支的实际源码}]{../supporting_materials/code/run_pre_freeze_scenarios.py}

\lstinputlisting[style=paperpython,firstline=305,lastline=357,firstnumber=305,caption={全期背景下限与候选选择的实际源码}]{../supporting_materials/code/run_pre_freeze_scenarios.py}

\section{模型验证程序}

验证程序从导出的任务表重新建立区间和容量占用，不调用主调度器的插入函数。除逐任务检查释放、截止、准备和设备能力外，还独立检查设备与患者区间重叠、跨室转运以及每个5分钟槽的医生并发。下面直接摘录\path{verify_unified_schedule.py}中的互斥与容量复核；问题二、三精确子问题分别见两个CP-SAT程序。

\lstinputlisting[style=paperpython,basicstyle=\ttfamily\fontsize{8.5pt}{9.5pt}\selectfont,firstline=220,lastline=364,firstnumber=220,caption={互斥、转运与医生容量复核的实际源码}]{../supporting_materials/code/verify_unified_schedule.py}
